\documentclass[aspectratio=169,10pt]{beamer}
% !TeX program = pdflatex
\usepackage[utf8]{inputenc}
\usepackage[T5]{fontenc}
\usepackage[vietnamese]{babel}
\usepackage{lmodern}
\usepackage{tikz}
\usetikzlibrary{arrows.meta,positioning,shapes.geometric,calc}
\usepackage{booktabs,tabularx,array}
\usepackage{amsmath,amssymb}
\usepackage{ragged2e}

\definecolor{Navy}{HTML}{3435B4}
\definecolor{Teal}{HTML}{4144C5}
\definecolor{Mint}{HTML}{ECECFA}
\definecolor{Gold}{HTML}{7378D2}
\definecolor{BarLR}{HTML}{087E8B}
\definecolor{BarRL}{HTML}{D28B34}
\definecolor{Coral}{HTML}{B44343}
\definecolor{MaxFill}{HTML}{DCEBF3}
\definecolor{MaxLine}{HTML}{7F9EAF}
\definecolor{MinFill}{HTML}{F5DEDA}
\definecolor{MinLine}{HTML}{BF8882}
\definecolor{Ink}{HTML}{233445}
\definecolor{Pale}{HTML}{F4F7F9}
\definecolor{Muted}{HTML}{718096}
\usetheme{Madrid}
% Madrid is a built-in Beamer lecture theme: rounded title/block boxes and an infolines footer.
\setbeamertemplate{navigation symbols}{}
\setlength{\parskip}{3pt}

\tikzset{
  decisionnode/.style={regular polygon,regular polygon sides=3,text=Ink,
    font=\fontsize{4.8}{5.8}\selectfont\bfseries,minimum size=6.5mm,
    text width=3.1mm,text height=1.3mm,text depth=0pt,align=center,inner sep=.2pt},
  maxnode/.style={decisionnode,shape border rotate=0,draw=MaxLine,fill=MaxFill},
  minnode/.style={decisionnode,shape border rotate=180,draw=MinLine,fill=MinFill},
  leaf/.style={rounded corners=2pt,draw=Muted!45,fill=Pale,text=Ink,font=\small,minimum width=9mm,minimum height=6mm},
  cutleaf/.style={leaf,draw=Coral,fill=Coral!10,text=Coral},
  edge/.style={draw=Muted!70,line width=.8pt},
  goodedge/.style={draw=Teal,line width=1.5pt},
  cutedge/.style={draw=Coral,dashed,line width=1.1pt}
}
\newcommand{\pill}[2]{\tikz[baseline=-.55ex]\node[rounded corners=2pt,fill=#1!13,text=#1,inner xsep=5pt,inner ysep=2pt,font=\scriptsize\bfseries]{#2};}
\newcommand{\smallsource}[1]{\vfill{\scriptsize\textcolor{Muted}{Nguồn: #1}}}

\title[Minimax và Alpha--Beta]{Cây trò chơi: Minimax và Alpha--Beta}
\author[Nhóm 10]{Nhóm 10\\{\small Hoàng Trung Hiếu, Trần Nguyên Hiển Minh,\\ Lê Bảo Ngân, Nguyễn Hoàng Anh}}
\institute[]{Biểu diễn tri thức và tìm kiếm nâng cao\\ IAI -- UET -- VNU}
\date{}

\begin{document}

% 1 / 9 — Title slide
\begin{frame}
\titlepage
\end{frame}

% 2 / 9
\begin{frame}{Cây T0 trong đề bài}
\centering
\small Mỗi nút trong chọn một trong hai con; MAX tối đa, MIN tối thiểu utility. Lá là trạng thái kết thúc, utility nhìn từ MAX.
\par\vspace{2pt}
\begin{tikzpicture}[x=1cm,y=1cm,baseline=(current bounding box.north)]
\tikzset{
  tmax/.style={maxnode},
  tmin/.style={minnode}
}
\node[tmax] (R) at (0,0) {R};
\node[tmin] (A) at (-3.6,-.9) {A};
\node[tmin] (B) at (3.6,-.9) {B};
\node[tmax] (A1) at (-5.4,-1.8) {A1};
\node[tmax] (A2) at (-1.8,-1.8) {A2};
\node[tmax] (B1) at (1.8,-1.8) {B1};
\node[tmax] (B2) at (5.4,-1.8) {B2};
\node[tmin] (A11) at (-6.3,-2.7) {A11};
\node[tmin] (A12) at (-4.5,-2.7) {A12};
\node[tmin] (A21) at (-2.7,-2.7) {A21};
\node[tmin] (A22) at (-.9,-2.7) {A22};
\node[tmin] (B11) at (.9,-2.7) {B11};
\node[tmin] (B12) at (2.7,-2.7) {B12};
\node[tmin] (B21) at (4.5,-2.7) {B21};
\node[tmin] (B22) at (6.3,-2.7) {B22};
\foreach \p/\c in {R/A,R/B,A/A1,A/A2,B/B1,B/B2,
  A1/A11,A1/A12,A2/A21,A2/A22,B1/B11,B1/B12,B2/B21,B2/B22}
  \draw[edge] (\p)--(\c);
\foreach \x/\v/\id in {
  -6.75/9/L01,-5.85/12/L02,-4.95/11/L03,-4.05/10/L04,
  -3.15/8/L05,-2.25/15/L06,-1.35/14/L07,-.45/13/L08,
  .45/4/L09,1.35/5/L10,2.25/3/L11,3.15/6/L12,
  4.05/100/L13,4.95/101/L14,5.85/90/L15,6.75/95/L16}
  {
  \node[leaf,minimum width=8.5mm,minimum height=1mm,font=\scriptsize\bfseries] (\id) at (\x,-3.6) {\v};
  \node[font=\tiny,text=Muted] at (\x,-3.9) {\id};
  }
\foreach \p/\c in {
  A11/L01,A11/L02,A12/L03,A12/L04,A21/L05,A21/L06,A22/L07,A22/L08,
  B11/L09,B11/L10,B12/L11,B12/L12,B21/L13,B21/L14,B22/L15,B22/L16}
  \draw[edge] (\p)--(\c);
\end{tikzpicture}
\smallsource{Đề bài, tr. 2. Các nhãn L01--L16 và vị trí lá được giữ nguyên.}
\end{frame}

% 3 / 9
\begin{frame}{Cây kiểm thử N1 và N2}
\begin{center}\small Seed 2579; các số ở lá là utility, theo thứ tự trái sang phải.\end{center}
\vspace{2pt}
\begin{columns}[T,onlytextwidth]
\column{.54\textwidth}
\centering
\textbf{N1: gốc MAX, 16 lá}\par\vspace{3pt}
\begin{tikzpicture}[x=.42cm,y=.78cm]
\tikzset{
compactmax/.style={
  maxnode,minimum size=4.5mm,text width=2.4mm,text height=1.1mm,
  font={\fontsize{3.7}{4.4}\selectfont\bfseries}
},
compactmin/.style={
  minnode,minimum size=4.5mm,text width=2.4mm,text height=1.1mm,
  font={\fontsize{3.7}{4.4}\selectfont\bfseries}
},
utility/.style={
  leaf,
  minimum width=3.5mm,
  minimum height=4.5mm,
  inner sep=.2pt,
  font=\tiny\bfseries
}
}
\node[compactmax] (n1r) at (0,0) {R};
\node[compactmin] (n1a) at (-3.6,-1) {A};
\node[compactmin] (n1b) at (3.6,-1) {B};
\node[compactmax] (n1a1) at (-5.4,-2) {A1};
\node[compactmax] (n1a2) at (-1.8,-2) {A2};
\node[compactmax] (n1b1) at (1.8,-2) {B1};
\node[compactmax] (n1b2) at (5.4,-2) {B2};
\node[compactmin] (n1a11) at (-6.3,-3) {A11};
\node[compactmin] (n1a12) at (-4.5,-3) {A12};
\node[compactmin] (n1a21) at (-2.7,-3) {A21};
\node[compactmin] (n1a22) at (-.9,-3) {A22};
\node[compactmin] (n1b11) at (.9,-3) {B11};
\node[compactmin] (n1b12) at (2.7,-3) {B12};
\node[compactmin] (n1b21) at (4.5,-3) {B21};
\node[compactmin] (n1b22) at (6.3,-3) {B22};
\foreach \p/\c in {n1r/n1a,n1r/n1b,n1a/n1a1,n1a/n1a2,n1b/n1b1,n1b/n1b2,
  n1a1/n1a11,n1a1/n1a12,n1a2/n1a21,n1a2/n1a22,
  n1b1/n1b11,n1b1/n1b12,n1b2/n1b21,n1b2/n1b22}
  \draw[edge] (\p)--(\c);
\foreach \x/\v/\i in {
  -6.75/24/01,-5.85/22/02,-4.95/3/03,-4.05/-3/04,
  -3.15/21/05,-2.25/13/06,-1.35/18/07,-.45/-3/08,
  .45/-5/09,1.35/29/10,2.25/-5/11,3.15/21/12,
  4.05/-9/13,4.95/29/14,5.85/29/15,6.75/-3/16}
  \node[utility] (n1l\i) at (\x,-3.85) {\v};
\foreach \p/\i in {n1a11/01,n1a11/02,n1a12/03,n1a12/04,
  n1a21/05,n1a21/06,n1a22/07,n1a22/08,
  n1b11/09,n1b11/10,n1b12/11,n1b12/12,
  n1b21/13,n1b21/14,n1b22/15,n1b22/16}
  \draw[edge] (\p)--(n1l\i);
\end{tikzpicture}

\column{.43\textwidth}
\centering
\textbf{N2: gốc MIN, 8 lá}\par\vspace{3pt}
\begin{tikzpicture}[x=.60cm,y=.78cm]
\tikzset{
  compactmax/.style={maxnode,minimum size=4.5mm,text width=2.4mm,text height=1.1mm,font=\fontsize{3.7}{4.4}\selectfont\bfseries},
  compactmin/.style={minnode,minimum size=4.5mm,text width=2.4mm,text height=1.1mm,font=\fontsize{3.7}{4.4}\selectfont\bfseries},
  utility/.style={
  leaf,
  minimum width=6mm,
  minimum height=4mm,
  inner sep=.2pt,
  font=\tiny\bfseries
}
}
\node[compactmin] (n2r) at (0,0) {R};
\node[compactmax] (n2a) at (-2.4,-1) {A};
\node[compactmax] (n2b) at (2.4,-1) {B};
\node[compactmin] (n2a1) at (-3.6,-2) {A1};
\node[compactmin] (n2a2) at (-1.2,-2) {A2};
\node[compactmin] (n2b1) at (1.2,-2) {B1};
\node[compactmin] (n2b2) at (3.6,-2) {B2};
\foreach \p/\c in {n2r/n2a,n2r/n2b,n2a/n2a1,n2a/n2a2,n2b/n2b1,n2b/n2b2}
  \draw[edge] (\p)--(\c);
\foreach \x/\v/\i in {-4.2/14/01,-3/-15/02,-1.8/11/03,-.6/12/04,
  .6/-19/05,1.8/-6/06,3/17/07,4.2/-13/08}
  \node[utility] (n2l\i) at (\x,-2.85) {\v};
\foreach \p/\i in {n2a1/01,n2a1/02,n2a2/03,n2a2/04,
  n2b1/05,n2b1/06,n2b2/07,n2b2/08}
  \draw[edge] (\p)--(n2l\i);
\end{tikzpicture}
\end{columns}
\smallsource{du\_lieu\_nhom.json, seed 2579. N1 có 16 utility; N2 có 8 utility.}
\end{frame}

% 4 / 9
\begin{frame}{Minimax: chọn kết quả có thể bảo đảm}
\begin{columns}[T,onlytextwidth]
\column{.51\textwidth}
\textbf{Tính từ lá lên gốc}
\vspace{2pt}
\begin{align*}
V(s)&=u(s) && \text{nếu $s$ là lá}\\[-2pt]
V(s)&=\max_{c\in C(s)} V(c) && \text{nếu $s$ là MAX}\\[-2pt]
V(s)&=\min_{c\in C(s)} V(c) && \text{nếu $s$ là MIN}
\end{align*}
\begin{beamercolorbox}[rounded=true,sep=7pt,wd=\linewidth]{block body}
\textbf{T0, thứ tự LR} \,\,\textcolor{Muted}{(trái sang phải)}\\
\small $A=\min(10,13)=10$; \quad $B=\min(4,100)=4$\\
\textbf{$R=\max(10,4)=10$ \;→\; chọn A}
\end{beamercolorbox}
\vspace{3pt}\small Đường chơi tối ưu:\\
$R\to A\to A1\to A12\to L04$.
\par\scriptsize Khi hòa, giữ con xét trước: MAX dùng $>$, MIN dùng $<$.

\column{.46\textwidth}
\centering
\textbf{T0: giá trị truyền từ dưới lên}\par\vspace{7pt}
\begin{tikzpicture}[x=1cm,y=1cm,
  valuemax/.style={maxnode},
  valuemin/.style={minnode}]
\node[valuemax] (r) at (0,0) {R};
\node[valuemin] (a) at (-1.5,-1.25) {A};
\node[valuemin] (b) at (1.5,-1.25) {B};
\node[valuemax] (a1) at (-2.2,-2.5) {A1};
\node[valuemax] (a2) at (-.75,-2.5) {A2};
\node[valuemax] (b1) at (.75,-2.5) {B1};
\node[valuemax] (b2) at (2.2,-2.5) {B2};
\draw[goodedge] (r)--(a);\draw[edge] (r)--(b);
\draw[goodedge] (a)--(a1);\draw[edge] (a)--(a2);
\draw[edge] (b)--(b1);\draw[edge] (b)--(b2);
\foreach \n/\v in {r/10,a/10,b/4,a1/10,a2/13,b1/4,b2/100}
  \node[below=1pt of \n,text=Ink,font=\small\bfseries] {\v};
\end{tikzpicture}
\pill{Teal}{Đọc đủ 16 lá}\par\vspace{4pt}
\small B có lá 101, nhưng MIN chọn B1 để ép điểm xuống 4.

\end{columns}
\note{Chính sách MAX: chọn A; nếu MIN chọn A1 thì đi A12, nếu chọn A2 thì đi A22. Các nút A1, A2, B1, B2 đều là MAX; các tầng sâu hơn được lược trong hình.}
\end{frame}

% 5 / 9
\begin{frame}{Alpha--Beta: bỏ nhánh không thể đổi lựa chọn}
\begin{columns}[T,onlytextwidth]
\column{.35\textwidth}
\textbf{Đã có 10 từ A.}\par
\begin{itemize}\setlength{\itemsep}{5pt}
  \item Giữ $\alpha=10$ khi sang B.
  \item B11 đọc 4, cắt L10; B12 đọc 3, cắt L12.
  \item B trả cận trên 4: không thể tốt hơn A, nên bỏ B2.
\end{itemize}
\vspace{4pt}
\small $\alpha$: mức điểm MAX đã có thể bảo đảm từ các lựa chọn đã xét.\par\vspace{3pt}
$\beta$: mức điểm MIN đã có thể giới hạn MAX nhận từ các lựa chọn đã xét.\par\vspace{4pt}
\scriptsize MAX cắt khi $best\ge\beta$ (LOWER); MIN khi $best\le\alpha$ (UPPER).

\column{.61\textwidth}
\centering
\begin{tikzpicture}[x=1cm,y=.9cm,
  pmax/.style={maxnode},
  pmin/.style={minnode}]
\node[pmax] (r) at (0,0) {R};
\node[pmin] (a) at (-2,-1) {A};
\node[pmin] (b) at (1.3,-1) {B};
\node[pmax] (b1) at (.1,-2) {B1};
\node[pmax,opacity=.35] (b2) at (3,-2) {B2};
\node[pmin] (b11) at (-.65,-3) {B11};
\node[pmin] (b12) at (.95,-3) {B12};
\node[leaf] (l09) at (-1.1,-4) {4};
\node[cutleaf,opacity=.5] (l10) at (-.15,-4) {5};
\node[leaf] (l11) at (.8,-4) {3};
\node[cutleaf,opacity=.5] (l12) at (1.75,-4) {6};
\node[rounded corners=2pt,draw=Coral,dashed,fill=Coral!4,
  text=Muted,font=\scriptsize,align=center,inner sep=5pt] (skipped) at (3,-3.45)
  {B21, B22\\L13--L16\\chưa duyệt};
\draw[goodedge] (r)--(a);\draw[edge] (r)--(b);
\draw[edge] (b)--(b1);\draw[cutedge] (b)--(b2);
\draw[edge] (b1)--(b11);\draw[edge] (b1)--(b12);
\draw[edge] (b11)--(l09);\draw[cutedge] (b11)--(l10);
\draw[edge] (b12)--(l11);\draw[cutedge] (b12)--(l12);
\draw[cutedge,opacity=.4] (b2)--(skipped);
\node[text=Coral,font=\small\bfseries] at (2.15,-1.5) {$\times$};
\node[text=Coral,font=\small\bfseries] at (-.4,-3.5) {$\times$};
\node[text=Coral,font=\small\bfseries] at (1.35,-3.5) {$\times$};
\node[text=Teal,font=\small\bfseries,align=center] at (-2,-2.05) {10\\đã xét};
\node[text=Muted,font=\scriptsize] at (.9,.05) {$\alpha=10$};
\foreach \n/\label in {l09/L09,l10/L10,l11/L11,l12/L12}
  \node[below=1pt of \n,font=\tiny,text=Muted] {\label};
\end{tikzpicture}
\par\scriptsize T0 rút gọn: cây con A đã xét; nét đỏ chỉ phần bị bỏ.
\end{columns}
\vspace{5pt}
\begin{beamercolorbox}[rounded=true,sep=5pt,wd=\linewidth]{block body}
\centering\textbf{Cùng kết quả: 10 / A}\quad 16 lá $\to$ 10 lá tiết kiệm 37,5\%\\
\scriptsize $22=12+10$ nút thăm; 3 lần cắt / 3 cạnh bỏ / 6 lá chưa đọc.
\end{beamercolorbox}
\note{Alpha và beta là các ngưỡng truyền theo đường tìm kiếm, có thể đến từ lựa chọn tại nút tổ tiên, không mặc nhiên là cận của giá trị thật tại nút đang xét. Ví dụ B nhận alpha=10 từ lựa chọn A ở gốc, dù giá trị thật của B là 4. Tại B11, cửa sổ vào là (10,+inf). Đọc L09=4, MIN cắt L10 và trả UPPER: V(B11)<=4. Không reset alpha khi chuyển từ A sang B.}
\end{frame}

% 6 / 9
\begin{frame}{Đúng điểm vẫn có thể chọn sai nước đi}
\begin{columns}[T,onlytextwidth]
\column{.50\textwidth}
\centering
\textbf{Cây T1}\par\vspace{5pt}
\begin{tikzpicture}[x=1.05cm,y=.88cm]
\node[maxnode] (r) at (0,0) {R};
\node[minnode] (a) at (-1.45,-1) {A};
\node[minnode] (b) at (1.45,-1) {B};
\draw[goodedge](r)--(a);\draw[edge](r)--(b);
\node[leaf] (10a) at (-2,-2) {10};\node[leaf] (12) at (-0.9,-2) {12};
\node[leaf] (10b) at (0.9,-2) {10};\node[cutleaf] (2) at (2,-2) {2};
\draw[edge](a)--(10a);\draw[edge](a)--(12);\draw[edge](b)--(10b);\draw[cutedge](b)--(2);
\node[font=\scriptsize,text=Coral] at (1.9,-2.55) {chưa đọc ở LR};
\node[font=\scriptsize,text=Muted] at (0,-3) {Giá trị thật: A=10, B=2, R=10};
\end{tikzpicture}
\column{.47\textwidth}
\vspace{2pt}
\begin{itemize}\setlength{\itemsep}{8pt}
  \item Duyệt LR: A cho MAX giá trị 10; sang B thì $\alpha=10$.
  \item MIN tại B đọc lá 10 rồi cắt lá 2; trả về \textbf{10, UPPER}.
  \item UPPER chỉ nói $V(B)\le10$; đó không nhất thiết là giá trị thật.
\end{itemize}
\begin{beamercolorbox}[rounded=true,sep=8pt,wd=\linewidth]{block body}
Nếu chọn nhánh bằng so sánh $\ge$, kết quả có thể là $(10,B)$ dù MAX thực nhận 2.\\[2pt]
\textbf{MAX chỉ đổi action khi $score>best$.}\\
\small Điều kiện cắt vẫn cho phép dấu bằng.
\end{beamercolorbox}
\end{columns}
\end{frame}

% 7 / 9
\begin{frame}{Kiểm chứng cả giá trị lẫn hành động}
\centering
\begin{tikzpicture}[node distance=0.45cm]
\node[rounded corners=3pt,fill=Navy,text=white,minimum width=2.6cm,minimum height=.72cm,font=\small\bfseries] (n1) {Tính tay độc lập};
\node[rounded corners=3pt,fill=Navy,text=white,minimum width=2.35cm,minimum height=.72cm,font=\small\bfseries,right=of n1] (n2) {Chạy chương trình};
\node[rounded corners=3pt,fill=Navy,text=white,minimum width=2.25cm,minimum height=.72cm,font=\small\bfseries,right=of n2] (n3) {Đối chiếu trace};
\node[rounded corners=3pt,fill=Navy,text=white,minimum width=2.3cm,minimum height=.72cm,font=\small\bfseries,right=of n3] (n4) {Kiểm thử tự động};
\draw[-{Stealth},thick,Muted] (n1)--(n2);\draw[-{Stealth},thick,Muted] (n2)--(n3);\draw[-{Stealth},thick,Muted] (n3)--(n4);
\end{tikzpicture}
\vspace{.3cm}
\begin{columns}[T,onlytextwidth]
\column{.57\textwidth}
\begin{itemize}\setlength{\itemsep}{7pt}
  \item T0: Minimax và Alpha--Beta cùng $(10,A)$; trace khớp nhánh cắt tính tay.
  \item T1: kiểm cả value và action để bắt lỗi đúng điểm, sai nhánh.
  \item Bao phủ điểm âm, utility 0, gốc MIN, hòa LR/RL, số đếm và trace.
\end{itemize}
\column{.40\textwidth}
\small
\begin{tabularx}{\linewidth}{>{\RaggedRight\arraybackslash}X>{\RaggedRight\arraybackslash}X}
\toprule\textbf{Thử lỗi} & \textbf{Test phát hiện}\\\midrule
Khởi tạo MAX bằng 0 & Test 1: dữ liệu toàn âm\\[5pt]
MAX đổi action bằng $\ge$ & Test 3: T1 trả $(10,B)$ thay vì $(10,A)$\\
\bottomrule
\end{tabularx}
\end{columns}
\vfill
\begin{center}\pill{Teal}{9/9 test PASS}\quad\small\textcolor{Muted}{Báo cáo ghi nhận 9 test nhóm, không có lỗi.}\end{center}
\smallsource{BaoCao.pdf, tr. 8--9; test\_nhom.py; output/ket\_qua\_test.txt}
\end{frame}

% 8 / 9
\begin{frame}{Thứ tự duyệt đổi mức tiết kiệm}
\small 24 cấu hình bắt buộc + 8 cấu hình riêng; N1/N2 dùng seed 2579.
\par\vspace{5pt}
\begin{columns}[T,onlytextwidth]
\column{.58\textwidth}
\centering
\textbf{Số lá đã đọc}\par\vspace{6pt}
\begin{tikzpicture}[x=.5cm,y=.20cm]
\foreach \y in {0,4,8,12,16}{
  \draw[Muted!20] (-.35,\y)--(12.4,\y);
  \node[left,font=\scriptsize,text=Muted] at (-.35,\y) {\y};
}
\foreach \x/\lab/\mm/\lr/\rl in {0/T0/16/10/16,5/N1/16/9/16,10/N2/8/8/6}{
  \node[font=\scriptsize,text=Ink] at (\x+1.05,-1.1) {\lab};
  \fill[Navy] (\x,0) rectangle +(0.55,\mm);
  \fill[BarLR] (\x+.75,0) rectangle +(0.55,\lr);
  \fill[BarRL] (\x+1.5,0) rectangle +(0.55,\rl);
  \node[above,font=\scriptsize,text=Ink] at (\x+.28,\mm) {\mm};
  \node[above,font=\scriptsize,text=Ink] at (\x+1.03,\lr) {\lr};
  \node[above,font=\scriptsize,text=Ink] at (\x+1.78,\rl) {\rl};
}
\end{tikzpicture}
\vspace{3pt}
\par\centering\scriptsize
\textcolor{Navy}{\rule{5pt}{5pt} Minimax}\hspace{.18cm}%
\textcolor{BarLR}{\rule{5pt}{5pt} AB--LR}\hspace{.18cm}%
\textcolor{BarRL}{\rule{5pt}{5pt} AB--RL}
\column{.38\textwidth}
\centering\small
\begin{tabular}{l c}
\toprule\textbf{Cây} & \textbf{Value / action}\\\midrule
T0 & $10/A$\\
N1 (MAX) & $13/A$\\
N2 (MIN) & $-13/B$\\\bottomrule
\end{tabular}
\par\vspace{8pt}
\begin{beamercolorbox}[rounded=true,sep=8pt,wd=\linewidth]{block body}
\textbf{Không có thứ tự luôn tốt nhất}\\
LR tốt hơn trên T0/N1; RL tốt hơn trên N2.
\end{beamercolorbox}
\end{columns}
\vspace{7pt}\small
\textbf{Thay dữ liệu:} $L01=t\in\{10,11,14\}$ $\Rightarrow$ value $10,11,12$; AB/LR đọc $10,9,9$ lá.\\[3pt]
\textbf{Đổi thang:} $u\mapsto2u-7$ $\Rightarrow$ value $10\to13$, vẫn chọn A và đọc 10 lá.
\smallsource{BaoCao.pdf, tr. 7; output/bat\_buoc.csv; output/tu\_tao.csv; du\_lieu\_nhom.json}
\end{frame}

% 9 / 9
\begin{frame}{Bỏ bớt nhánh, vẫn cần bảo đảm chọn đúng}
\begin{columns}[T,onlytextwidth]
\column{.57\textwidth}
\begin{enumerate}\setlength{\itemsep}{8pt}
  \item Alpha--Beta giữ value/action ở gốc như Minimax với cùng thứ tự và quy tắc hòa.
  \item Thứ tự duyệt quyết định mức tiết kiệm; có trường hợp không giảm số lá.
  \item Kiểm cả điểm và hành động; phân biệt cận với giá trị thật.
\end{enumerate}
\column{.39\textwidth}
\begin{beamercolorbox}[rounded=true,sep=8pt,wd=\linewidth]{block body}
\textbf{Giới hạn khảo sát}\\
Cây tối đa 16 lá; chưa đo thời gian lặp hoặc khảo sát nhiều seed.\\[4pt]
\small Thời gian xấu nhất $O(b^d)$; bộ nhớ phụ bản cài đặt $O(N)$.\\[4pt]
\end{beamercolorbox}
\end{columns}
\vspace{9pt}
\begin{beamercolorbox}[rounded=true,sep=7pt,wd=\linewidth]{block body}
\textcolor{Teal}{\bfseries Demo T1 (3 phút):} dự đoán nhánh cắt, chỉ mã action/cận, rồi đổi LR/RL.\\[4pt]
\scriptsize Tái chạy:\\
\texttt{python run.py -{}-batch -{}-tag bat\_buoc}\\
\texttt{python -m unittest -v test\_nhom}
\end{beamercolorbox}
\note{Phân công: Hiếu phụ trách Minimax/B1, test 1--2; Minh phụ trách Alpha--Beta/B2--B4, test 3--4; Ngân phụ trách tham số/N1/N2, test 5--7; Hoàng Anh phụ trách thực nghiệm/tích hợp, test 8--9.}
\end{frame}

\end{document}
